\chapter{Présentation de l'entreprise d'accueil}

\section{Présentation détaillée de Prestige Maroc}
\textbf{Nom de l'entreprise :} Prestige Maroc \\
\textbf{Secteur d'activité :} Location de véhicules (Tourisme, Utilitaires, Longue durée) \\

Prestige Maroc est une agence locale spécialisée dans la location de véhicules pour les particuliers (B2C) et les entreprises (B2B). Fondée il y a près d'une décennie, l'entreprise s'est progressivement forgée une excellente réputation sur le marché national grâce à une stratégie axée sur le renouvellement constant de sa flotte et un service client de proximité. 

Ses missions principales s'articulent autour de trois axes stratégiques :
\begin{itemize}
    \item \textbf{La Mobilité Touristique :} Mise à disposition de véhicules de tourisme (citadines, SUV) pour une clientèle internationale et locale, particulièrement lors des hautes saisons estivales.
    \item \textbf{Le Leasing B2B :} Contrats de location longue durée (LLD) destinés aux entreprises locales souhaitant externaliser la gestion de leur parc automobile.
    \item \textbf{Le Service Premium :} Offre de transferts aéroportuaires et de location avec chauffeur privé pour une clientèle VIP.
\end{itemize}

\section{Environnement Socio-économique et Concurrence}
Le marché de la location de voitures au Maroc est hautement fragmenté. Prestige Maroc évolue dans un environnement concurrentiel dense, face à :
\begin{enumerate}
    \item \textbf{Les multinationales franchisées} (Hertz, Avis, Sixt) qui captent la majorité des flux aéroportuaires grâce à leurs systèmes de réservation mondiaux.
    \item \textbf{Les agences locales indépendantes} qui se battent principalement sur une guerre des prix, souvent au détriment de la qualité du service.
\end{enumerate}
Pour se démarquer de ces deux extrêmes, Prestige Maroc mise sur un positionnement intermédiaire : offrir la qualité de service d'une multinationale au prix d'une agence locale. Cependant, cette ambition est freinée par un manque cruel d'outils digitaux.

\section{Organigramme et Structuration Interne}
L'entreprise est structurée de manière agile avec plusieurs pôles interdépendants :
\begin{itemize}
    \item \textbf{Direction Générale :} Pilotage de la stratégie globale, gestion des partenariats B2B et validation des grilles tarifaires saisonnières.
    \item \textbf{Service Commercial :} Point de contact direct avec le client. Gère l'accueil physique, la négociation, la rédaction des contrats de location et la fidélisation.
    \item \textbf{Service Logistique (Parc automobile) :} Responsable du cycle de vie physique des véhicules. S'occupe de l'entretien mécanique régulier, du nettoyage, et surtout de la vérification méticuleuse de l'état des véhicules lors des phases de Check-in (retour) et Check-out (départ).
    \item \textbf{Service Comptabilité et Finance :} Assure le recouvrement des factures, le suivi des cautions bancaires, le paiement des fournisseurs (garages, assurances) et les déclarations fiscales mensuelles.
\end{itemize}

\section{Place du projet dans l'entreprise}
Mon stage a été réalisé en tant que Développeur Full-Stack et Chef de Projet technique, directement rattaché à la Direction Générale. Mon rôle ne consistait pas simplement à "créer un site web", mais à mener à bien une véritable \textbf{transformation digitale systémique}. L'objectif était d'auditer les processus manuels des différents services (Commercial, Logistique, Comptabilité) pour concevoir une solution logicielle sur-mesure (ERP/CRM) capable de centraliser, sécuriser et automatiser 100\% des flux de travail de l'entreprise.
